\section{Modelagem de banco de dados}
\label{sec:banco}
A modelagem distingue o esquema operacional da extensão relacional proposta. O banco da implementação Java possui somente a tabela \texttt{audit\_records}, criada pela migração Flyway \texttt{V1\_\_create\_audit\_records.sql}. O acesso usa Spring JDBC; não há mapeamento ORM nem rotina de adaptação de linhas históricas da versão Python. Marcos de execução são emitidos nos logs da aplicação, fora do banco operacional. A coleta estruturada de métricas e tentativas para o experimento ainda será desenvolvida. As tabelas adicionais desta seção permanecem propostas opcionais.

\subsection{Modelo conceitual (MER)}
A Figura~\ref{fig:mer-atual} apresenta o núcleo e a Figura~\ref{fig:mer-proposto} apresenta a extensão proposta. No núcleo atual, a entidade \textit{Registro de auditoria} representa o fato persistido e contém identificação, origem, instante, entidade, ator, conteúdo e hash. Não há relacionamento entre tabelas no esquema existente. O hash é atributo do registro, pois há exatamente um digest armazenado por linha; criar uma entidade \textit{EventHash} independente não acrescentaria informação necessária ao recorte e criaria uma sincronização adicional.

A extensão proposta introduz \textit{Log de processamento}, \textit{Tentativa de recuperação} e \textit{Amostra de métrica}. Um UUID pode aparecer em zero ou muitos logs, inclusive quando nenhuma linha de auditoria foi persistida. Cada log representa uma tentativa ou etapa observada. Um log pode originar zero ou muitas tentativas de recuperação; cada tentativa referencia um log. Métricas agregadas são associadas ao identificador da execução, e não obrigatoriamente a um único evento.

\capfig{mer-atual}{Modelo conceitual do núcleo existente}{fig:mer-atual}{0.25}
\capfig{mer-proposto}{Modelo conceitual da extensão de observabilidade proposta}{fig:mer-proposto}{0.40}

A associação entre registro e log é de correlação por UUID, não de existência obrigatória: um erro pode ocorrer antes da primeira inserção. Não se estabelece uma chave estrangeira obrigatória de log para \texttt{audit\_records}, pois isso impediria registrar justamente a falha de persistência. A tentativa de recuperação, por sua vez, referencia um log já existente. Essa distinção preserva a capacidade de registrar eventos rejeitados, falhas e reentregas sem pressupor sucesso.

\subsection{Modelo lógico}
No esquema atual, \texttt{audit\_records.event\_id} é chave primária; \texttt{entity\_id} possui o índice secundário \texttt{idx\_audit\_records\_entity\_id}; os atributos do evento e o hash são não nulos. A coluna \texttt{persisted\_at} é anulável no DDL, embora novas inserções recebam \texttt{clock\_timestamp()} por padrão. A Figura~\ref{fig:der-atual} representa somente a tabela implementada. Não existem chaves estrangeiras ou relacionamentos entre tabelas. A restrição de unicidade protege contra inserções concorrentes com o mesmo UUID. O hash é calculado na aplicação e armazenado como texto; não há restrição SQL de comprimento do digest.

\capfig{modelo-logico-atual}{Modelo lógico do banco implementado}{fig:der-atual}{0.38}

\subsection{Modelo físico implementado}
O arquivo \texttt{audit-storage/src/main/resources/db/migration/V1\_\_create\_audit\_records.sql} cria a tabela no PostgreSQL. Sua chave primária sobre \texttt{event\_id} é a única restrição de unicidade explícita. \texttt{entity\_id} recebe índice não único \texttt{idx\_audit\_records\_entity\_id}. Os nove campos obrigatórios são \texttt{event\_id}, \texttt{event\_type}, \texttt{entity\_type}, \texttt{entity\_id}, \texttt{actor\_id}, \texttt{source}, \texttt{occurred\_at}, \texttt{payload} e \texttt{content\_hash}. O décimo campo, \texttt{persisted\_at}, é anulável e possui \texttt{DEFAULT clock\_timestamp()}. Não existe chave estrangeira, gatilho, exclusão lógica ou tabela de usuários nesse esquema.

O Flyway aplica a migração V1 na inicialização da API síncrona e do consumidor sobre o banco novo desta versão. Não há migração automática do banco Python nem compatibilidade retroativa demonstrada. A garantia de idempotência combina a PK, o \texttt{INSERT ON CONFLICT DO NOTHING} e a comparação do \texttt{content\_hash} no \texttt{AuditRepository} após a tentativa de inserção. Portanto, o modelo físico deve ser entendido junto à regra implementada na aplicação, e não como uma restrição SQL de igualdade do conteúdo.

\subsection{Extensão lógica opcional}
A Figura~\ref{fig:der} apresenta, separadamente, a possível extensão de observabilidade. As entidades destacadas como \textit{propostas} não foram criadas pelo código e não integram o banco utilizado nos testes.

\capfig{modelo-logico}{Modelo lógico da extensão proposta, separado do esquema atual}{fig:der}{0.57}

Na extensão, \texttt{processing\_log.id}, \texttt{retry\_log.id} e \texttt{metrics.id} são chaves primárias propostas. \texttt{retry\_log.processing\_log\_id} referencia o log que motivou a tentativa. \texttt{run\_id} correlaciona os três conjuntos com o manifesto externo de execução. Não se cria uma tabela de execuções nesta proposta inicial; caso isso seja feito posteriormente, o manifesto poderá ser normalizado em uma entidade própria.

A unicidade proposta de \texttt{(run\_id, event\_id, component, attempt)} em logs identifica uma tentativa por componente. Índices por execução, UUID e instante apoiam conciliação. Status e unidades deverão ser validados pela camada de coleta. Migrações e consultas dessas tabelas são atividades futuras, não executadas para produzir este texto. A instrumentação deverá ser equivalente entre variantes e ter seu custo considerado na análise experimental.

\FloatBarrier
\subsection{Dicionário de dados}

\subsubsection{Tabela audit\_records --- existente}
A Tabela~\ref{tab:dic-audit} descreve o esquema efetivo. Os campos do evento e o hash são não nulos; \texttt{persisted\_at} admite nulo pelo DDL da migração Java. No recebimento HTTP, \texttt{event\_id}, \texttt{occurred\_at} e \texttt{payload} podem ser omitidos e recebem valores padrão.

{\small
\begin{longtable}{>{\raggedright\arraybackslash}p{0.23\textwidth}>{\raggedright\arraybackslash}p{0.22\textwidth}>{\raggedright\arraybackslash}p{0.46\textwidth}}
\caption{Dicionário do registro persistido}\label{tab:dic-audit}\\
\hline Campo & Tipo e tamanho & Obrigatoriedade, significado e validação\\ \hline
\endfirsthead
\hline Campo & Tipo e tamanho & Obrigatoriedade, significado e validação\\ \hline
\endhead
\texttt{event\_id} & UUID, 128 bits & PK; não nulo. UUID informado deve ser válido; uuid4 gerado se ausente. Deve ser preservado em reenvios experimentais.\\ \hline
\texttt{event\_type} & VARCHAR(100) & Não nulo; obrigatório na API, 1--100 caracteres. Tipo do fato observado; sem enumeração fechada implementada.\\ \hline
\texttt{entity\_type} & VARCHAR(100) & Não nulo; obrigatório, 1--100 caracteres. Tipo da entidade referida.\\ \hline
\texttt{entity\_id} & VARCHAR(200) & Não nulo; obrigatório, 1--200 caracteres; índice. Identificador da entidade e chave enviada ao Kafka.\\ \hline
\texttt{actor\_id} & VARCHAR(200) & Não nulo; obrigatório, 1--200 caracteres. Ator sintético relacionado ao fato.\\ \hline
\texttt{source} & VARCHAR(100) & Não nulo; obrigatório, 1--100 caracteres. Sistema ou componente de origem.\\ \hline
\texttt{occurred\_at} & TIMESTAMP WITH TIME ZONE & Não nulo. Instante de ocorrência; padrão UTC atual se omitido. Quando informado, o Java interpreta data ISO-8601 com offset. Não representa instante de commit.\\ \hline
\texttt{payload} & JSONB & Não nulo; objeto com padrão vazio. Sem limite de tamanho ou profundidade definido no DTO atual. Conteúdo contextual sintético.\\ \hline
\texttt{content\_hash} & TEXT & Não nulo. SHA-256 calculado, 64 caracteres hexadecimais; o tipo SQL não impõe esse comprimento. Inclui os campos completos do evento.\\ \hline
\texttt{persisted\_at} & TIMESTAMP WITH TIME ZONE & Anulável pelo DDL. Novas linhas recebem clock\_timestamp() durante INSERT. Não representa conclusão do commit.\\ \hline
\end{longtable}
\noindent{\footnotesize Fonte: Elaborado pelo autor (2026).}\par
}

\subsubsection{Tabela processing\_log --- proposta}
A Tabela~\ref{tab:dic-processing} descreve a proposta de uma tentativa por componente. Ela não substitui o evento original nem modifica sua identidade.

{\small
\begin{longtable}{>{\raggedright\arraybackslash}p{0.23\textwidth}>{\raggedright\arraybackslash}p{0.22\textwidth}>{\raggedright\arraybackslash}p{0.46\textwidth}}
\caption{Dicionário proposto dos logs de processamento}\label{tab:dic-processing}\\
\hline Campo & Tipo e tamanho & Obrigatoriedade e regra proposta\\ \hline
\endfirsthead
\hline Campo & Tipo e tamanho & Obrigatoriedade e regra proposta\\ \hline
\endhead
\texttt{id} & BIGINT identity & PK, obrigatório. Identifica o log.\\ \hline
\texttt{run\_id} & UUID & Obrigatório. Correlaciona com o manifesto de execução.\\ \hline
\texttt{event\_id} & UUID & Obrigatório para eventos identificáveis. Correlação, sem FK para o registro persistido. Erros sem UUID ficam no log técnico do coletor.\\ \hline
\texttt{component} & VARCHAR(30) & Obrigatório. Identifica API síncrona, API produtora ou consumidor.\\ \hline
\texttt{status} & VARCHAR(30) & Obrigatório. Vocabulário proposto: accepted, persisted, duplicate, rejected, failed ou conflict.\\ \hline
\texttt{attempt} & INTEGER & Obrigatório; maior ou igual a 1. Número da tentativa por evento e componente na execução.\\ \hline
\texttt{started\_at} & TIMESTAMPTZ & Obrigatório. Início da tentativa em UTC.\\ \hline
\texttt{finished\_at} & TIMESTAMPTZ & Opcional enquanto a tentativa estiver aberta; ao concluir, não anterior ao início.\\ \hline
\texttt{detail} & TEXT & Opcional. Diagnóstico sem credenciais ou conteúdo sensível.\\ \hline
\end{longtable}
\noindent{\footnotesize Fonte: Elaborado pelo autor (2026).}\par
}

\subsubsection{Tabela retry\_log --- proposta}
A Tabela~\ref{tab:dic-retry} descreve os campos da extensão para recuperação.
{\small
\begin{longtable}{>{\raggedright\arraybackslash}p{0.23\textwidth}>{\raggedright\arraybackslash}p{0.22\textwidth}>{\raggedright\arraybackslash}p{0.46\textwidth}}
\caption{Dicionário proposto das tentativas de recuperação}\label{tab:dic-retry}\\
\hline Campo & Tipo e tamanho & Obrigatoriedade e regra proposta\\ \hline
\endfirsthead
\hline Campo & Tipo e tamanho & Obrigatoriedade e regra proposta\\ \hline
\endhead
\texttt{id} & BIGINT identity & PK, obrigatório.\\ \hline
\texttt{processing\_log\_id} & BIGINT & FK obrigatória para processing\_log.id; identifica a tentativa anterior relacionada.\\ \hline
\texttt{reason} & VARCHAR(200) & Obrigatório. Motivo da retomada, sem dados sensíveis.\\ \hline
\texttt{retried\_at} & TIMESTAMPTZ & Obrigatório. Instante da nova tentativa; não apenas um agendamento futuro.\\ \hline
\end{longtable}
\noindent{\footnotesize Fonte: Elaborado pelo autor (2026).}\par
}
O vínculo proposto ao log fornece execução, UUID e componente sem replicar esses campos. A política automática de retry está implementada no consumidor e emite marcos JSON no fluxo de logs; sua existência não depende da criação desta tabela.

\subsubsection{Tabela metrics --- proposta}
A Tabela~\ref{tab:dic-metrics} descreve uma possível persistência relacional das métricas a serem coletadas no experimento Java.
{\small
\begin{longtable}{>{\raggedright\arraybackslash}p{0.23\textwidth}>{\raggedright\arraybackslash}p{0.22\textwidth}>{\raggedright\arraybackslash}p{0.46\textwidth}}
\caption{Dicionário proposto das amostras de métricas}\label{tab:dic-metrics}\\
\hline Campo & Tipo e tamanho & Obrigatoriedade e regra proposta\\ \hline
\endfirsthead
\hline Campo & Tipo e tamanho & Obrigatoriedade e regra proposta\\ \hline
\endhead
\texttt{id} & BIGINT identity & PK, obrigatório.\\ \hline
\texttt{run\_id} & UUID & Obrigatório. Correlaciona com o manifesto.\\ \hline
\texttt{variant} & VARCHAR(10) & Obrigatório: sync ou async.\\ \hline
\texttt{metric\_type} & VARCHAR(50) & Obrigatório. Identifica latência, throughput, erro, recurso ou backlog.\\ \hline
\texttt{value} & DOUBLE PRECISION & Obrigatório; valor finito. Ausência de amostra não deve ser registrada como zero.\\ \hline
\texttt{unit} & VARCHAR(20) & Obrigatório. Unidade explícita, como ms, eventos/s, bytes ou percentual.\\ \hline
\texttt{window\_start} & TIMESTAMPTZ & Obrigatório. Início da janela de medição.\\ \hline
\texttt{window\_end} & TIMESTAMPTZ & Obrigatório; igual ou posterior ao início.\\ \hline
\end{longtable}
\noindent{\footnotesize Fonte: Elaborado pelo autor (2026).}\par
}

A medição de conclusão utiliza o marco \texttt{commit\_completed} emitido após a transação. Nem \texttt{occurred\_at} nem \texttt{persisted\_at} substituem esse marco. Consultas periódicas oferecem uma observação posterior da linha, sujeita ao intervalo de consulta, à duração da requisição e ao escalonamento; portanto, o intervalo configurado sozinho não estabelece um limite rígido para todo o atraso de observação.

\textbf{Síntese.} O esquema existente garante unicidade e conserva conteúdo e hash. O modelo proposto amplia a rastreabilidade de falhas e métricas sem impor uma FK que impeça registrar tentativas fracassadas e sem apresentar essas tabelas como já implementadas.
